\section{求导穿过积分号需要控制差商}
\label{sec:12-Real-Analysis-Support/04-Limits-and-Integration/03}

你要优化一个期望损失 \(L(\theta)=\E_{X\sim P}[\ell(\theta,X)]\)。代码里的做法是采一批样本，对每个样本算 \(\nabla_\theta\ell(\theta,x_i)\)，再平均。这算的是梯度的期望。你想要的是期望的梯度。这两个量凭什么相等？

自动微分不会替你回答。它忠实地对采样程序做链式法则，得到的永远是逐样本求导再平均；两者相等与否是一条分析结论，需要单独验证，而验证的对象不是被积函数本身。

把问题写出来。对
\[
F(t)=\int_E f(t,x)\,d\mu(x),
\]
求导意味着
\[
\frac{F(t+h)-F(t)}h=\int_E\frac{f(t+h,x)-f(t,x)}h\,d\mu(x),
\]
再让 \(h\to0\)。左边的极限是 \(F'(t)\)，右边要把极限推进积分号——被积的那一列函数是差商，不是 \(f\)，也不是 \(\partial_t f\)。于是上一节的控制收敛定理照样适用，只是需要控制的对象换成了差商。严格地说，导数是连续极限，用定理时要对任意趋于零的序列 \(h_k\) 分别应用。

\subsection{控制函数值管不住导数}

为什么必须盯住差商，一个例子说得很清楚。取 \(f_n(x)=\frac{\sin(nx)}n\)，它一致收敛到零，函数值被 \(1/n\) 均匀压住，压得要多小有多小。可 \(f_n'(x)=\cos(nx)\)，振幅始终是 1，随 \(n\) 增大来回摆得越来越快，根本不收敛。

\begin{center}
\begin{tikzpicture}[>=stealth]
\draw[->] (0,0) -- (5.9,0) node[right] {\footnotesize \(x\)};
\draw[domain=0:5.2,samples=300] plot ({\x},{0.2*sin(720*\x)});
\draw[very thick] (-0.4,-0.75) -- (0.4,0.75);
\node[right] at (0.35,0.95) {\footnotesize 切线依旧很陡};
\node[right] at (3.3,0.5) {\footnotesize \(f_n=\sin(nx)/n\)};
\node at (2.6,-1.7) {\footnotesize 函数值被压在零附近};
\begin{scope}[shift={(8.4,0)}]
\draw[->] (0,0) -- (5.9,0) node[right] {\footnotesize \(x\)};
\draw[domain=0:5.2,samples=300] plot ({\x},{1.15*cos(720*\x)});
\node[right] at (3.3,1.5) {\footnotesize \(f_n'=\cos(nx)\)};
\node at (2.6,-1.7) {\footnotesize 导数振幅纹丝不动};
\end{scope}
\end{tikzpicture}
\end{center}

\emph{曲线可以被压得任意扁，斜率却完全不受这份压制约束。}

左图那条几乎贴着横轴的曲线，切线斜率接近 \(\pm1\)。把 \(n\) 加倍，曲线更扁，切线还是那么陡。函数值的一致上界与导数的大小之间没有任何蕴含关系，所以任何只约束 \(\abs{f_n}\) 的条件都不足以交换求导与极限。

对应的正确条件是把控制加在导数一侧：若 \(f_n\) 在 \([a,b]\) 上可微，存在一点 \(x_0\) 使 \((f_n(x_0))\) 收敛，且 \(f_n'\) 一致收敛到 \(g\)，则 \(f_n\) 一致收敛到某个 \(f\) 且 \(f'=g\)。证明走微积分基本定理，把 \(f_n(x)=f_n(x_0)+\int_{x_0}^xf_n'\) 中的积分与极限交换，用的是导数列的一致收敛。同一条原则换个包装还会再出现：幂级数在收敛半径内部可以逐项求导，因为在任意更小的闭区间上导数级数一致收敛；到了收敛边界，这份控制失效，必须单独讨论。

\subsection{一个够用的充分条件与它的对账}

回到含参积分。设在 \(t_0\) 的某个邻域 \(I\) 内，对几乎所有 \(x\)，\(\partial_t f(t,x)\) 存在；存在可积函数 \(g\) 使 \(\abs{\partial_t f(t,x)}\le g(x)\) 对\emph{所有} \(t\in I\) 成立；并且至少有一个参数点上 \(f(t,x)\) 可积、可测性成立。那么 \(F\) 在 \(t_0\) 可导且
\[
F'(t_0)=\int_E\partial_t f(t_0,x)\,d\mu(x).
\]

推导只有一步转换：由一元中值定理，对每个 \(x\)，差商等于某个中间点 \(\xi\) 处的 \(\partial_t f(\xi,x)\)，而 \(\xi\) 落在 \(t_0\) 与 \(t_0+h\) 之间，因此被 \(g(x)\) 控制。差商族有了与 \(h\) 无关的可积包络，控制收敛定理接手。

条件对账里最要紧的是包络的作用范围。只知道 \(\partial_t f(t_0,x)\) 可积是不够的——中值点 \(\xi\) 随 \(h\) 在邻域里游走，你不知道它落在哪，所以 \(g\) 必须罩住整个邻域内的所有 \(t\)。可微性提供差商的逐点极限，中值定理把差商翻译成导数，邻域一致的可积包络提供控制，某点可积保证 \(F\) 本身有定义而不是恒为无穷。

一个直接例子：\(F(t)=\int_0^1e^{-tx}\,dx\)，\(t\ge0\)。此时 \(\partial_t e^{-tx}=-xe^{-tx}\)，而对 \(t\ge0\)、\(x\in[0,1]\) 有 \(\abs{xe^{-tx}}\le x\)，\(x\) 在 \([0,1]\) 上可积且与 \(t\) 无关。于是 \(F'(t)=-\int_0^1xe^{-tx}\,dx\)。这个例子有闭式 \(F(t)=(1-e^{-t})/t\) 可以对照，但求导穿过积分号的做法不依赖闭式，这正是它在没有解析积分的模型上仍然可用的原因。

\subsection{积分限也动时，边界项不能漏}

若上下限本身依赖参数，
\[
F(t)=\int_{a(t)}^{b(t)}f(t,x)\,dx,
\]
在相应的连续与可微条件下，
\[
\begin{aligned}
F'(t)&=f(t,b(t))\,b'(t)-f(t,a(t))\,a'(t)\\
&\quad+\int_{a(t)}^{b(t)}\partial_t f(t,x)\,dx.
\end{aligned}
\]
前两项来自区域端点的移动，最后一项来自被积函数自身的变化。

统计里最容易踩的坑是分布的支撑随参数移动。\(\mathrm{Uniform}(0,\theta)\) 的密度是 \(\frac1\theta\mathbf1_{[0,\theta]}(x)\)，对 \(\theta\) 求导时若只对 \(\frac1\theta\) 动手而忽略上限 \(\theta\) 的移动，得到的``score''会漏掉边界贡献，进而让基于它的信息量计算和渐近结论一起失真。这也是这一族分布在极大似然理论中屡屡被当作反例的原因。

\subsection{对期望求梯度靠的是同一条规则}

回到开头的问题。若数据分布 \(P\) 不依赖 \(\theta\)，只要 \(\nabla_\theta\ell(\theta,x)\) 在参数邻域内有共同的可积包络，就可以写
\[
\nabla_\theta\,\E_{X\sim P}[\ell(\theta,X)]=\E_{X\sim P}[\nabla_\theta\ell(\theta,X)],
\]
逐样本求导再平均因此是合法的。重参数化梯度做的正是把参数从分布挪到被积函数里：把 \(z\sim q_\theta\) 写成 \(z=T_\theta(\epsilon)\)、\(\epsilon\sim q_0\) 之后，随机性所在的测度不再含 \(\theta\)，剩下的就是这一条规则的直接应用，它在 \hyperref[sec:14-Deep-Learning/06-Variational-Autoencoders/01]{``VAE 把不可积后验变成可训练下界''}中被反复使用。

若测度本身带参数，导数还会落到密度上：
\[
\nabla_\theta\,\E_{P_\theta}[f(X)]=\E_{P_\theta}\bigl[f(X)\,\nabla_\theta\log p_\theta(X)\bigr],
\]
这是 score function 估计量的核心恒等式，成立条件是支撑不随 \(\theta\) 变动加上前述的可积控制，进一步的用法见 \hyperref[sec:11-Mathematical-Statistics/03-Sufficient-Statistics-and-Information/02]{``Score 与 Fisher 信息衡量局部可辨认性''}。支撑会移动时，上一小节的边界项就要补回来。

对于似然函数也是同一套：写下 \(\E_\theta[\nabla_\theta\log p_\theta(X)]=0\) 这条基本等式时，用到的就是求导与积分的交换，正则条件不满足的模型里它不成立。

三节走到这里，被交换的对象从函数换成差商，条件的形式没有变过：找一个与极限参数无关的可积量，把整族摁住。写下交换那一步之前，值得先说出这个量是谁。
